%! The Statistical Cycle and Types of Data — Unit 1, Lesson 1.1 — Guided Notes
\documentclass[10pt]{article}
\usepackage{saar-article}
\usepackage{saar-boxes}
\newcommand{\TallMath}[1]{$\displaystyle #1\rule[-1.4em]{0pt}{3.2em}$}

\begin{document}

\pageheader{Unit 1, Lesson 1.1}{Guided Notes}
% NAMESTRIP: no \namedateperiod here. The student writes their name once, on the
% cover this component is stapled behind. See references/conventions.md.

\vspace{0.15in}

\begin{objectivebox}
\small
By the end of this lesson, I will be able to\dots
\begin{enumerate}[leftmargin=1.5em, itemsep=2pt, topsep=2pt]
  \item name the four stages of the \textbf{data cycle} and say what each stage
        produces for the next one;
  \item write a question that data can actually answer --- one whose answers
        \emph{vary}, that names the \emph{group}, and that says what will be
        \emph{measured};
  \item decide from the question alone whether the data I need are
        \textbf{categorical} or \textbf{quantitative};
  \item organize collected data into a frequency table and turn counts into
        percents;
  \item write a conclusion that names the group it applies to --- and say what
        the data do \emph{not} tell me.
\end{enumerate}
\end{objectivebox}

\boxguard
\begin{vocabbox}
\small
Fill in each term as we build it together.
\termblanklong{Data cycle}
\termblanklong{Statistical question}
\termblanklong{Variable of interest}
\termblanklong{Frequency table}
\termblanklong{Relative frequency}
\termblanklong{Conclusion in context}
\end{vocabbox}

\boxguard[12]
\begin{hookbox}
\small
The Riverbend Student Council has one hour on the district's agenda. They want to
ask for \textbf{more bike racks} outside the main entrance.

\begin{center}
\textit{``We think a lot of students would bike if there were somewhere to park.''}
\end{center}

The district will not act on ``we think.'' \textbf{What would the council have to
do first --- and in what order?} Write the first two steps you would take.
\writelines{2}
\end{hookbox}

\boxguard
\begin{notesbox}{1. The Four Stages of the Data Cycle}
\small
Every study in this course runs the same four stages, in order --- and then
starts over. Fill in the name of each stage as we go.

\begin{center}
\renewcommand{\arraystretch}{1.35}
\begin{tabularx}{0.97\linewidth}{@{} c p{4.3cm} X @{}}
\toprule
\textbf{Stage} & \textbf{Name} & \textbf{What it hands to the next stage} \\
\midrule
1 & \blank{3.8cm} & A question that names a group and says what will be measured. \\
2 & \blank{3.8cm} & The raw values --- one row per individual. \\
3 & \blank{3.8cm} & A table or display that makes the pattern visible. \\
4 & \blank{3.8cm} & A conclusion, stated in context, plus the next question. \\
\bottomrule
\end{tabularx}
\end{center}

\vspace{3pt}
\textbf{Why the order matters:} each stage can only use what the stage before it
handed over. If stage 1 never says \emph{who} the study is about, then stage 4
cannot say who the conclusion applies to.

\vspace{2pt}
And it is a \textbf{cycle}, not a list: the answer at stage 4 almost always
raises a new question, which sends you back to stage \blank{1.2cm}.
\end{notesbox}

\boxguard
\begin{notesbox}{2. Stage 1 --- A Question Data Can Answer}
\small
A \textbf{statistical question} has to pass three tests.

\begin{center}
\renewcommand{\arraystretch}{1.3}
\begin{tabularx}{0.95\linewidth}{@{} c X @{}}
\toprule
\textbf{Test} & \textbf{Ask yourself} \\
\midrule
1 & Do the answers \blank{3cm} from one individual to the next? \\
2 & Does the question name the \blank{3cm} it is about? \\
3 & Does it say what will be \blank{3cm} or recorded? \\
\bottomrule
\end{tabularx}
\end{center}

\vspace{4pt}
Each question below fails at least one test. Say which test it fails, then repair
it.

\begin{center}
\renewcommand{\arraystretch}{1.5}
\begin{tabularx}{0.97\linewidth}{@{} p{4.4cm} p{1.5cm} X @{}}
\toprule
\textbf{Weak question} & \textbf{Fails} & \textbf{Repaired question} \\
\midrule
How do I get to school?          & \blank{1cm} & \blank{6cm} \\
Do people like biking?           & \blank{1cm} & \blank{6cm} \\
How far away is Riverbend?       & \blank{1cm} & \blank{6cm} \\
\bottomrule
\end{tabularx}
\end{center}

\vspace{2pt}
\textit{The council's real question: ``How do Riverbend students usually get to
school?'' Answers vary, the group is named, and what gets recorded is clear.}
\end{notesbox}

\boxguard
\begin{notesbox}{3. The Question Decides the Type of Data}
\small
You do not wait until the data arrive to find out what type they are. The
\textbf{question} tells you --- and the type tells you how you will be allowed to
summarize the answers.

\begin{center}
\renewcommand{\arraystretch}{1.5}
\begin{tabularx}{0.98\linewidth}{@{} X p{3.1cm} p{2.2cm} p{2.4cm} @{}}
\toprule
\textbf{Question} & \textbf{What you record} & \textbf{Type} &
\textbf{Summarize with} \\
\midrule
How do Riverbend students usually get to school?
  & \blank{2.7cm} & \blank{1.8cm} & \blank{2cm} \\
How many minutes is a Riverbend student's one-way trip?
  & \blank{2.7cm} & \blank{1.8cm} & \blank{2cm} \\
How old are the bikes students ride to school?
  & \blank{2.7cm} & \blank{1.8cm} & \blank{2cm} \\
\bottomrule
\end{tabularx}
\end{center}

\vspace{3pt}
The thing you record is the \textbf{variable of interest}. Naming it out loud is
the fastest way to catch a question that is not ready:
if you cannot finish the sentence \emph{``for each student I will
record \underline{\hspace{2.6cm}},''} the question is not finished either.
\end{notesbox}

\boxguard[30]
\begin{notesbox}{4. Stages 2 and 3 --- Collect, Then Organize}
\small
The council asked $60$ Riverbend students how they usually get to school. That is
stage 2. Now organize the $60$ answers into a \textbf{frequency table} --- counts
first, then each count as a percent of the total.

\begin{center}
\renewcommand{\arraystretch}{1.4}
\begin{tabularx}{0.9\linewidth}{@{} l c X @{}}
\toprule
\textbf{How they get to school} & \textbf{Frequency (count)} &
\textbf{Relative frequency (percent)} \\
\midrule
Bus   & 27 & $27/60 = 0.45 = 45\%$ \\
Car   & 15 & \blank{4.5cm} \\
Walk  & 12 & \blank{4.5cm} \\
Bike  &  6 & \blank{4.5cm} \\
\midrule
\textbf{Total} & \blank{1.4cm} & \blank{4.5cm} \\
\bottomrule
\end{tabularx}
\end{center}

\vspace{3pt}
Show the work for the \textbf{Bike} row:

\begin{work}
  \dfrac{6}{60} &= 0.10 \\
                &= 10\%
\end{work}

\vspace{2pt}
\textbf{Check:} the four percents must add to \blank{1.6cm}. If they do not, a
count is wrong or a student was left out.

\vspace{2pt}
\textit{Notice what stage 3 did: the counts were already there at stage 2, but
until they were organized, nobody could see that biking is the smallest group.}
\end{notesbox}

\boxguard
\begin{notesbox}{5. Stage 4 --- Analyze and Communicate}
\small
A conclusion is not a number. It is a sentence that says \textbf{what the data
show}, \textbf{who it applies to}, and \textbf{in what units}.

\vspace{3pt}
Scale the sample percent up to the whole school ($900$ students):

\begin{work}
  0.10 \times 900 &= 90 \text{ students}
\end{work}

\vspace{3pt}
Fill in the council's conclusion:

\vspace{2pt}
\textit{``Of the \blank{1.4cm} students we surveyed, \blank{1.4cm}\% bike to
school. If our sample is representative, that is about \blank{1.4cm} of
Riverbend's \blank{1.4cm} students.''}

\vspace{4pt}
\textbf{Now the hard part --- what the data do \emph{not} say.} The council wants
to claim students \emph{would} bike if there were racks. Does this survey support
that claim? \blank{2cm}

Why not? \writeline
\end{notesbox}

\boxguard
\begin{notesbox}{6. The Cycle Turns}
\small
The council's answer produced a \emph{new} question --- ``How many students would
bike if there were racks?'' --- and no data on hand can answer it. So they go back
to stage \blank{1.2cm} and run the cycle again with a new question.

\vspace{3pt}
Each scenario below breaks down at one stage. Write the stage number.

\begin{enumerate}[leftmargin=1.6em, itemsep=4pt, topsep=3pt]
  \item The council only surveyed students already in the bike rack area.
        \hfill \blank{1.2cm}
  \item The report says ``most students want racks,'' but nobody was ever asked
        about racks. \hfill \blank{1.2cm}
  \item The council asked ``How do you get to school?'' without saying which
        students they meant. \hfill \blank{1.2cm}
  \item The $60$ answers were left as a list of $60$ words, so nobody could see
        which group was largest. \hfill \blank{1.2cm}
\end{enumerate}
\end{notesbox}

\boxguard
\begin{practicebox}
\small
Six of the surveyed students also reported their one-way trip time, in minutes:
$25, \; 10, \; 5, \; 20, \; 12, \; 18$.

\textbf{(a)} Find the mean trip time.

\begin{work}
  \bar{x} &= \dfrac{25+10+5+20+12+18}{6} \\
          &= \dfrac{90}{6} \\
          &= 15 \text{ minutes}
\end{work}

\textbf{(b)} What percent of the $60$ students surveyed walk to school?

\begin{work}
  \dfrac{12}{60} &= 0.20 \\
                 &= 20\%
\end{work}

\textbf{(c)} Predict how many of Riverbend's $900$ students walk to school.

\begin{work}
  0.20 \times 900 &= 180 \text{ students}
\end{work}

\textbf{(d)} Write the stage-4 conclusion for part (a). Name the group, the
number, and the units.
\writeline
\end{practicebox}

\end{document}
